
\chapter{مسیر یادگیری Frontend از صفر}

\section{چرا مستقیم وارد React نمی‌شویم؟}

React روی مفاهیم Browser و JavaScript بنا شده است. اگر DOM، Event، Fetch و State را نفهمیم، React تبدیل به مجموعه‌ای از دستورهای حفظی می‌شود.

بنابراین یادگیری Frontend در این پروژه مرحله‌ای است.

\section{مرحله F1: HTML}

موضوعات:
\begin{itemize}
  \item document structure؛
  \item semantic elements؛
  \item heading hierarchy؛
  \item link vs button؛
  \item form؛
  \item label/input؛
  \item nav/main/article/section؛
  \item accessibility basics.
\end{itemize}

تمرین پروژه:
\begin{itemize}
  \item Idea Card استاتیک؛
  \item Login form؛
  \item Navigation ساده.
\end{itemize}

\begin{learningbox}
اول باید بفهمی چرا \eng{button} برای Action و \eng{a} برای Navigation استفاده می‌شود. این تفاوت بعداً روی Accessibility و Keyboard behavior تأثیر دارد.
\end{learningbox}

\section{مرحله F2: CSS}

موضوعات:
\begin{itemize}
  \item selector؛
  \item cascade؛
  \item specificity؛
  \item box model؛
  \item width/height؛
  \item margin/padding؛
  \item Flexbox؛
  \item Grid؛
  \item responsive؛
  \item media query؛
  \item CSS variable؛
  \item logical properties؛
  \item RTL/LTR.
\end{itemize}

تمرین:
\begin{itemize}
  \item بازسازی IdeaCard پروتوتایپ؛
  \item ساخت Bottom Navigation؛
  \item responsive کردن Feed.
\end{itemize}

\section{مرحله F3: Browser JavaScript}

با وجود تجربه Node.js، باید Browser-specific concepts را یاد بگیری.

\subsection{DOM}

\begin{latin}
\begin{lstlisting}
const button = document.querySelector("#save-button");
button.addEventListener("click", () => {
  console.log("saved");
});
\end{lstlisting}
\end{latin}

موضوعات:
\begin{itemize}
  \item DOM tree؛
  \item querySelector؛
  \item event؛
  \item event target؛
  \item form submit؛
  \item preventDefault.
\end{itemize}

\subsection{Fetch}

\begin{latin}
\begin{lstlisting}
const response = await fetch("/api/v1/ideas/");
const data = await response.json();
\end{lstlisting}
\end{latin}

باید بفهمی:
\begin{itemize}
  \item Promise؛
  \item async/await؛
  \item network failure؛
  \item HTTP error با exception جاوااسکریپت یکی نیست؛
  \item JSON parsing.
\end{itemize}

\section{مرحله F4: npm و package.json}

مفاهیم:
\begin{itemize}
  \item package؛
  \item dependency؛
  \item devDependency؛
  \item script؛
  \item lockfile؛
  \item semantic versioning در حد عملی.
\end{itemize}

\section{مرحله F5: TypeScript}

مثال JavaScript:
\begin{latin}
\begin{lstlisting}
function formatIdea(idea) {
  return idea.title;
}
\end{lstlisting}
\end{latin}

TypeScript:
\begin{latin}
\begin{lstlisting}
type Idea = {
  id: number;
  title: string;
  body: string;
};

function formatIdea(idea: Idea): string {
  return idea.title;
}
\end{lstlisting}
\end{latin}

موضوعات:
\begin{itemize}
  \item primitive types؛
  \item array/object؛
  \item type alias؛
  \item interface؛
  \item union؛
  \item optional field؛
  \item function type؛
  \item generic در حد لازم؛
  \item narrowing؛
  \item unknown vs any.
\end{itemize}

\begin{warningbox}
برای سریع شدن، \eng{any} را به Default تبدیل نکن. اگر داده API را نمی‌دانیم، Contract را تعریف می‌کنیم.
\end{warningbox}

\section{مرحله F6: React Fundamentals}

\subsection{Component}

\begin{latin}
\begin{lstlisting}
function IdeaCard() {
  return <article>...</article>;
}
\end{lstlisting}
\end{latin}

\subsection{Props}

\begin{latin}
\begin{lstlisting}
type IdeaCardProps = {
  title: string;
  readingTime: string;
};

function IdeaCard({ title, readingTime }: IdeaCardProps) {
  return (
    <article>
      <h2>{title}</h2>
      <span>{readingTime}</span>
    </article>
  );
}
\end{lstlisting}
\end{latin}

\subsection{State}

مثلاً local state برای باز/بسته بودن Menu.

\subsection{Event}

Save click، Language switch، Form submit.

\subsection{List Rendering}

Feed:
\begin{latin}
\begin{lstlisting}
{ideas.map((idea) => (
  <IdeaCard key={idea.id} idea={idea} />
))}
\end{lstlisting}
\end{latin}

\section{مرحله F7: Hooks}

به ترتیب:
\begin{itemize}
  \item \tech{useState}
  \item \tech{useEffect}
  \item \tech{useMemo} فقط وقتی دلیل وجود دارد
  \item custom hook
  \item context.
\end{itemize}

\begin{warningbox}
useEffect محل انجام هر نوع Logic نیست. قبل از استفاده باید بفهمیم آیا واقعاً Side Effect داریم یا می‌توان مقدار را مستقیم derive کرد.
\end{warningbox}

\section{مرحله F8: Routing}

Routeهای اصلی:
\begin{latin}
\begin{lstlisting}
/login
/register
/onboarding
/feed
/ideas/:ideaId
/explore
/library
/creators/:creatorId
/notifications
/profile
\end{lstlisting}
\end{latin}

Messaging later:
\begin{latin}
\begin{lstlisting}
/messages
/messages/:conversationId
\end{lstlisting}
\end{latin}

\section{مرحله F9: Server State}

بعد از اینکه fetch دستی را فهمیدی، TanStack Query وارد می‌شود.

مفاهیم:
\begin{itemize}
  \item queryKey؛
  \item queryFn؛
  \item stale/fresh؛
  \item cache؛
  \item mutation؛
  \item invalidation؛
  \item optimistic update؛
  \item infinite query.
\end{itemize}

\section{تمرین‌های Frontend در ترتیب واقعی}

\begin{enumerate}
  \item Static IdeaCard با HTML/CSS؛
  \item Responsive mobile feed؛
  \item language direction toggle؛
  \item fetch ساده از endpoint آزمایشی؛
  \item TypeScript type برای Idea؛
  \item React IdeaCard؛
  \item React feed از mock array؛
  \item Routing؛
  \item fetch real API؛
  \item TanStack Query؛
  \item Save mutation؛
  \item optimistic UI.
\end{enumerate}

\section{معیار فهم Frontend}

باید بتوانی توضیح دهی:
\begin{itemize}
  \item DOM چیست؟
  \item React چه مشکلی را نسبت به DOM manipulation دستی حل می‌کند؟
  \item Props و State چه فرقی دارند؟
  \item چرا Server State را با TanStack Query نگه می‌داریم؟
  \item TypeScript در Runtime وجود دارد یا Compile-time؟
  \item Vite چه کاری می‌کند و چه کاری نمی‌کند؟
\end{itemize}
